%!TEX root = ../main.tex

\chapter{LLM-Only Baseline}\label{ch:llm_baseline}

Before introducing retrieval, a baseline is established that relies solely on the
\gls{llm}'s parametric knowledge to answer questions. This baseline quantifies the
contribution of the training-time knowledge of the model and provides a lower bound
for the \gls{rag} experiments that follow.


% ===========================================================================
\section{System Design}

\subsection{Language Model}

All \gls{llm} calls in this work are routed through an Azure~OpenAI endpoint
using the \texttt{gpt-5.4-nano} model series. The model receives each prompt as a
single user turn and returns a text completion, using the deployment's default
sampling parameters; reproducibility of reported results is ensured by the response
cache described below. The \gls{api} client is implemented in
\texttt{src/llm/client.py} and is shared across all pipelines; no second inference
client is introduced at any point.

To avoid redundant \gls{api} calls during iterative development, all prompt-response
pairs are cached on disk keyed by the SHA-256 hash of the prompt string, in
\texttt{experiments/cache/}. Generation across large example sets is parallelised
using a thread pool (\texttt{ThreadPoolExecutor}) with exponential back-off on
rate-limit errors.

\subsection{Prompt Design}

The \gls{llm}-only prompt instructs the model to act as a government-services agent
and to answer in one to two plain sentences. The exact abstention phrase is fixed
so that the \texttt{is\_abstention()} evaluation function can detect it reliably.
The prompt used for this baseline is:

\begin{quote}
\small\ttfamily
You are a government services agent. Answer the customer's question in 1-2 plain
sentences. No bullet points, no bold text, no preamble. If you don't know, reply
exactly: ``Unfortunately, no relevant information is found.''

Customer: \{question\}\\
Agent:
\end{quote}

No retrieved context is provided. The model must answer from parametric knowledge
alone, with the dialogue history available only implicitly through the question text
(which may be decontextualised in later turns).


% ===========================================================================
\section{Results}

\subsection{Initial 50-Example Evaluation}

A preliminary evaluation on a balanced set of 50 examples (drawn before the dev-500
subset was formalised) was used to validate the evaluation pipeline and the prompt
design. Results are shown in Table~\ref{tab:llm:initial}; earlier prompt versions and
their results are documented in Appendix~\ref{ax:log:pilot}.

\begin{table}[ht]
    \centering
    \caption{LLM-only baseline on the initial 50-example set.}
    \label{tab:llm:initial}
    \begin{tabular}{lrr}
        \toprule
        Metric & Value \\
        \midrule
        Exact Match (EM)    & 0.0\% \\
        Token \fone{}        & 14.2\% \\
        \fone{} (answerable) & 15.0\% \\
        \fone{} (unanswerable) & 1.1\% \\
        Abstain rate         & 0.0\% \\
        \bottomrule
    \end{tabular}
\end{table}

The model does not abstain on any question in the absence of a retrieval signal,
meaning it always produces an answer even for questions that have no supporting
document. The \fone{} on unanswerable examples (1.1\%) is near zero because the
model generates plausible-sounding but unsupported responses that share little
token overlap with the reference abstention phrase.

\subsection{Dev-500 Evaluation}

On the fixed 500-example dev subset, the \gls{llm}-only system achieves a token
\fone{} of $16.3\%$ with a 95\% bootstrap confidence interval of
$[15.0, 17.6]$ (Table~\ref{tab:llm:dev500}). Abstain recall is low (19.4\%),
confirming that without a retrieval signal the model rarely produces the expected
abstention phrase. The false-abstention rate of 7.7\% indicates that even this
low recall comes at the cost of some incorrect abstentions on answerable questions.

\begin{table}[ht]
    \centering
    \caption{LLM-only baseline on the dev-500 subset ($n = 500$, 95\% bootstrap CI).}
    \label{tab:llm:dev500}
    \begin{tabular}{lr}
        \toprule
        Metric & Value \\
        \midrule
        Token \fone{} [CI]   & 16.3\% $[15.0, 17.6]$ \\
        Abstain recall [CI]  & 19.4\% $[6.5, 35.5]$ \\
        False abstention     & 7.7\% \\
        Abstain precision    & 14.3\% \\
        \bottomrule
    \end{tabular}
\end{table}


% ===========================================================================
\section{Analysis}

The 16.3\% \fone{} establishes that the model possesses some relevant knowledge
about government services from its pre-training data, but is far from reliable for
this task. The low abstain recall (19.4\%) confirms that without a retrieval signal
providing an empty or irrelevant context, the model will almost always attempt an
answer. This is consistent with findings from the broader \gls{rag} literature:
parametric models tend to confabulate rather than abstain when queried outside
their confident knowledge range.

The relatively high confidence-interval width for abstain recall (CI width 29.0~pp)
reflects the small number of unanswerable examples in the dev-500 subset ($n = 31$),
which makes abstention metrics noisy. For this reason, paired bootstrap tests on
abstention metrics are treated as illustrative rather than definitive.

This baseline motivates the \gls{rag} pipeline developed in the next chapter: the
retrieval step is expected both to improve answer quality on answerable questions
(by providing relevant context) and to enable more reliable abstention (by providing
an empty or irrelevant context on unanswerable inputs, which the generation prompt
can then treat as a signal to abstain).
